\chapter*{循迹之外：电磁学中的电像法}
\addcontentsline{toc}{chapter}{循迹之外：电磁学中的电像法}

\let\oldthesection\thesection
\let\oldthesubsection\thesubsection

\setcounter{section}{0}
\renewcommand{\thesection}{\arabic{section}}
\renewcommand{\thesubsection}{\thesection.\arabic{subsection}}
\setcounter{MgCtr}{0}


\section{问题的引入——导体与点电荷}

在GPA中\mgnote{见第8章导体与静电平衡，特别是\refleaftext{law8.3}。}，我们已经探讨了导体在电场中的性质。下面我们不妨更进一步，讨论如何解决点电荷与导体并存时的那类问题。

显然，若一个电荷系统的电荷分布已经给定，那么空间中任一位置的场强原则上均可以求得——只要利用点电荷产生的场强公式再加上叠加原理就能办到。但当空间存在导体时，问题就变得复杂多了。因为此时我们已知的是一组"等势面"，希望求得的结果是一组电荷分布（包括自由电荷和感应电荷），这组电荷分布形成的场强与点电荷的场强叠加后，最终要满足导体表面电势处处相等的结果。

为了稍微量化这其中的困难，我们不妨从线性代数的视角审视。若将空间中所有可能的电荷分布（包括自由电荷和感应电荷）视为一个向量空间 $\mathcal{V}_{\rho}$，将空间中各点的电势函数视为另一个向量空间 $\mathcal{V}_{\phi}$，则无导体时的场强计算对应于一个明确的正向线性算子 $\hat{L}$：给定电荷分布 $\rho$，通过库仑积分直接求得电势分布 $\phi$。然而，当存在导体时，情况便发生了翻转——我们已知的是导体表面电势满足一个常数向量 $C$（即等势面），而未知的是表面上的感应面电荷密度 $\sigma$。这等同于要求解一个算子方程：
\[\hat{K}(\sigma) = C - \phi_{\text{ext}},\]
其中 $\hat{K}$ 是一个定义在导体表面上的积分算子\mgnote{将 $\sigma$ 映射为电势的线性算子：$(\hat{K}\sigma)(\vec{r})=\frac{1}{4\pi\varepsilon_0}\int_S\frac{\sigma(\vec{r}\,')}{|\vec{r}-\vec{r}\,'|}\dif S'$}。

这件事之所以困难，在于两层障碍。第一，积分算子 $\hat{K}$ 本身是\itr{ill-posed}{不适定}\mgnote{违反 Hadamard 适定性条件；此处破坏的是稳定性——小扰动导致解的大幅震荡}的——微小的电荷扰动就可能引起电势分布的剧烈变化，这使得直接反演极其敏感。第二，该映射不具有局域性——表面上某一点的电荷变化，会通过长程的库仑相互作用影响到表面所有点的电势，导致离散化后的系数矩阵是\textbf{稠密矩阵}（每个元素都非零，$N$ 单元需 $O(N^2)$ 存储）而非稀疏矩阵。若将导体表面剖分为 $N$ 个单元，直接求解这个线性方程组需要 $O(N^3)$ 的计算量——当 $N$ 稍大时，这就是一场计算灾难。

而幸运的是，对于某些形状特殊的导体，电荷分布会最终导向人脑也能轻易计算的特殊解。找到这组特殊解的方法，便是这次循迹之外的主题——\itr{method of images}{电像法}。下面我们举一个例子，让读者直观感受电像法的作用。

\begin{wrapfigure}{r}{0.33\textwidth}
    \centering
    \thisincludegraphics[width=\linewidth]{electroimage_4.png}
    \caption{一对等量异号点电荷的电场线（实线）和等势面（虚线）。}
\end{wrapfigure}

设有一对点电荷，带电量分别为 $+q$、$-q$，相距 $2a$。我们知道，这一对电荷在空间中产生的电场线和等势面是容易画出的。现在考察其中一个等势面——假想我们有一个形状恰好与这个等势面相吻合的薄金属片，把它准确地安置在这个等势面上，同时调整此金属片的电势恰好为该等势面的电势值，那么整个空间的电场将不会发生任何改变。

上述分析事实上已经解决了一个新问题：一块给定电势的曲面导体（其形状与一对点电荷的某个等势面相同），在此导体外放一个点电荷 $+q$（位置与图中相同），那么导体表面外、与点电荷同侧空间中的场强分布，与一对点电荷（$+q$ 和 $-q$）所激发产生的场强分布完全相同。在求解中，我们不需要处理导体表面分布复杂的感应电荷，而只需去寻找一个假想的点电荷（称为\itr{image charge}{像电荷}）$-q$，便能轻松解决问题。

然而，如此巧妙的简化有一个前置条件：这组解必须是\textbf{唯一的}。如果同一组边界条件可以对应多个不同的电场分布，那么我们找到的"像电荷解"就只是众多可能中的一个，并没有真正解决问题。因此，在正式进入电像法之前，我们必须先推导一个前置结论——\itr{uniqueness theorem}{唯一性定理}。


\section{唯一性定理——物理视角}

\subsection{边值问题的提出}

在进入定理之前，我们需要先明确自己要解决的是什么问题。

实际中提出的静电学问题，大多不是已知电荷分布求电场分布，而是通过一定的电极来控制或实现某种电场分布。这里问题的出发点（已知的前提），除给定各带电导体的几何形状、相互位置外，往往再给定下列条件之一：
\begin{Itemize}
    \item 每个导体的电势 $U_k$；
    \item 每个导体上的总电量 $Q_k$。
\end{Itemize}
其中 $k = 1, 2, \ldots$ 为导体的编号。寻求的答案是上述条件（称为\itr{boundary conditions}{边界条件}）下电场的恒定分布。这类问题称为静电场的边值问题。

这里不谈边值问题如何解决，而我们要问一个更基本的问题：给定一组边界条件，空间能否存在不同的恒定电场分布？唯一性定理对此的回答是干脆的——\textbf{不能}。换句话说，边界条件可将空间里电场的恒定分布唯一地确定下来。

\subsection{三个引理}

在证明唯一性定理之前，我们先做点准备工作。为简单起见，暂把研究的问题限定为一组导体，除此之外的空间里没有电荷。

\begin{law}[引理一：无电荷空间电势无极值]
    在无电荷的空间里，电势不可能有极大值和极小值。
\end{law}

用反证法。设电势 $U$ 在空间某点 $P$ 极大，则在 $P$ 点周围的所有邻近点上，梯度 $\nabla U$ 必都指向 $P$ 点，即场强 $\vec{E} = -\nabla U$ 的方向都是背离 $P$ 点的（见图1a）。这时我们作一个很小的闭合面 $S$ 把 $P$ 点包围起来，穿过 $S$ 的电通量为
\[\Phi_E = \oint_{(S)} \vec{E} \cdot \dif \vec{S} > 0.\]
根据高斯定理\mgnote{\refleaftext{law8.1}}，$S$ 面内必然包含正电荷。然而这违背了"空间无电荷"的前提。因此 $U$ 不可能有极大值。

用同样的方法可以证明，$U$ 不可能有极小值（参见图2b）——若 $P$ 点为极小值，则所有电场线都指向 $P$ 点，电通量为负，面内必有负电荷，同样违背前提。

\begin{singlefigure}[引理一的证明——电势极值不存在]{electroimage_1.png}[0.7]
    （a）若 $P$ 点为极大值，电场线全部背离 $P$ 点，闭合面内必有正电荷；
    （b）若 $P$ 点为极小值，电场线全部指向 $P$ 点，闭合面内必有负电荷。
\end{singlefigure}

\begin{law}[引理二：零电势导体的空间]
    若所有导体的电势为 $0$，则导体以外空间的电势处处为 $0$。
\end{law}

这个引理的证明几乎就是引理一的直接推论。因为电势在无电荷空间里的分布是连续变化的，若空间有电势大于 $0$（或小于 $0$）的点，而边界上又处处等于 $0$，在空间必出现电势的极大（或极小）值，这直接违背引理一。

不难看出，本引理可稍加推广：若在完全由导体所包围的空间里各导体的电势都相等（设为 $U_0$），则空间电势等于常量 $U_0$。

\begin{law}[引理三：不带电导体的电势]
    若所有导体都不带电，则各导体的电势都相等。
\end{law}

用反证法。设各导体电势不全相等，则其中必有一个电势最高的，设它是导体 $1$。如图2所示，电场线只可能从导体 $1$ 出发到达其余导体 $2$、$3$、$\ldots$，而不可能反过来——毕竟电场线永远从高电势走向低电势。于是我们得到这样的结论：导体 $1$ 的表面上任何地方都只能是电场线的起点，不可能是终点，即此导体表面只有正电荷而无负电荷，从而它带的总电量不可能为 $0$。这违背了前提。

将引理三与引理二结合起来，可以进一步推论：在所有导体都不带电的情况下，空间各处的电势也和导体一样，等于同一常量。

\begin{singlefigure}[引理三的证明——电场线方向]{electroimage_2.png}[0.55]
    电场线只能从电势最高的导体出发，到达电势较低的导体。若所有导体不带电，则不存在电势差。
\end{singlefigure}

\subsection{叠加原理}

电场是服从叠加原理的——场强服从矢量叠加，电势服从代数叠加。由这个朴素的原理出发，我们可以得到一条对唯一性定理至关重要的推论。

在给定各带电导体的几何形状、相互位置后，赋予它们两组边界条件：
\begin{Itemize}
    \item 给定每个导体的电势为 $U_{\text{I}k}$（或总电量为 $Q_{\text{I}k}$）；
    \item 给定每个导体的电势为 $U_{\text{II}k}$（或总电量为 $Q_{\text{II}k}$）。
\end{Itemize}
设 $U_{\text{I}}$、$U_{\text{II}}$ 分别是满足边界条件（1）、（2）的恒定电势分布，则它们的线性组合 $U = aU_{\text{I}} + bU_{\text{II}}$ 必定是满足下列边界条件的恒定分布：
\begin{Itemize}
    \item 给定每个导体的电势为 $U_k = aU_{\text{I}k} + bU_{\text{II}k}$（或总电荷 $Q_k = aQ_{\text{I}k} + bQ_{\text{II}k}$）。
\end{Itemize}

作为一个特例，取 $U_{\text{I}k} = U_{\text{II}k}$（或 $Q_{\text{I}k} = Q_{\text{II}k}$）和 $a=1,\; b=-1$，则 $U = U_{\text{I}} - U_{\text{II}}$ 是满足"所有导体电势为 $0$（或总电量为 $0$）"的恒定分布。这意味着：\textbf{两个具有相同边界条件的解之差，满足零边界条件。}这一洞察是证明唯一性定理的最后一块拼图。

\subsection{唯一性定理}

\begin{law}[静电边值问题的唯一性定理]
    给定一组导体的几何形状和相互位置，以及每个导体上的电势 $U_k$（或总电量 $Q_k$），则空间中的恒定电场分布被唯一确定。
\end{law}

具体来说，分为两种情况：

\textbf{情况一：给定每个导体的电势。}
设对应同一组边值 $U_k \;(k=1,2,\ldots)$ 有两种恒定的电势分布 $U_{\text{I}}$ 和 $U_{\text{II}}$，则它们的差 $U = U_{\text{I}} - U_{\text{II}}$ 相当于所有导体上电势为 $0$ 时的恒定电势分布。根据引理二，空间电势 $U$ 恒等于 $0$，即 $U_{\text{I}}$ 恒等于 $U_{\text{II}}$，从而
\[\vec{E}_{\text{I}} = -\nabla U_{\text{I}} \quad \text{恒等于} \quad \vec{E}_{\text{II}} = -\nabla U_{\text{II}}.\]
电势分布和电场分布都是唯一的。

\textbf{情况二：给定每个导体上的总电量。}
第 $k$ 个导体上的总电量为
\[Q_k = \oint_{(S_k)} \sigma_e \dif S = \varepsilon_0 \oint_{(S_k)} E_n \dif S = -\varepsilon_0 \oint_{(S_k)} \frac{\partial U}{\partial n} \dif S,\]
式中 $S_k$ 为导体 $k$ 的表面，$\sigma_e$ 代表表面电荷密度，$n$ 代表法向。

设对应同一组边值 $Q_k \;(k=1, 2, \ldots)$ 有两种恒定电势分布 $U_{\text{I}}$ 和 $U_{\text{II}}$，即
\[-\varepsilon_0 \oint_{(S_k)} \frac{\partial U_{\text{I}}}{\partial n} \dif S = -\varepsilon_0 \oint_{(S_k)} \frac{\partial U_{\text{II}}}{\partial n} \dif S = Q_k, \quad (k = 1, 2, \ldots)\]
令 $U = U_{\text{I}} - U_{\text{II}}$，则
\[-\varepsilon_0 \oint_{(S_k)} \frac{\partial U}{\partial n} \dif S = 0, \quad (k = 1, 2, \ldots)\]
即 $U$ 相当于所有导体都不带电时的恒定电势分布。根据引理三后面的推论，在空间中有
\[U = U_{\text{I}} - U_{\text{II}} = \text{常量} \quad \text{或} \quad U_{\text{I}} = U_{\text{II}} + \text{常量},\]
此常量不影响其梯度：
\[\nabla U_{\text{I}} = \nabla U_{\text{II}},\]
即场强分布是完全一样的：
\[\vec{E}_{\text{I}} = \vec{E}_{\text{II}}.\]

电势中所差的常量与电势的参考点有关。只要各导体中有一个的电势确定了，其它导体以及空间的电势分布就唯一地确定下来了。

把上述证明推广到混合边界条件（即部分导体给定电势、部分给定总电量）的情形也是直接的，这里从略。

现在，我们不妨用刚学到的唯一性定理，来重新审视GPA中一个看似熟悉但并未严格证明的结论——\itr{electrostatic shielding}{静电屏蔽}。

\begin{ex}[静电屏蔽]
    取一任意形状的闭合金属壳，将它接地（见图3）。现从外面移来若干正或负的带电体。若腔内无带电体，则其中 $\vec{E}=0$（图3a）。反之，将带电体放进腔内，而壳外无带电体，则外部空间 $\vec{E}=0$（图3b）。今设想将a、b两图合并在一起（图3c），即壳外有与图a相同的带电体，腔内有与图b相同的带电体。现在我们要问：这时壳内、外电场的恒定分布是否仍分别与图a、b一样？
\end{ex}

\begin{so}[静电屏蔽]
    \begin{singlefigure}[静电屏蔽的三个情形]{electroimage_3.png}[0.85]
        （a）壳外有带电体、腔内无电荷时，腔内 $\vec{E}=0$；
        （b）腔内有带电体、壳外无电荷时，壳外 $\vec{E}=0$；
        （c）壳内外同时存在带电体时。
    \end{singlefigure}

    首先我们可以肯定，这是可能的。因为当外部电荷和电场分布如图a时，它在腔内不产生电场，从而腔内的带电体所处的环境和图b一样，故可产生与之相同的恒定分布。反之，当内部电荷和电场分布如图b时，它在壳外不产生电场，从而壳外带电体所处的环境和图a一样，故也可产生与之相同的恒定分布。以上的论述表明，壳内、外带电体同时存在时，若壳内、外的电荷和电场分别维持与a、b二图相同的分布，是可以达到静电平衡的。

    这里遗留的问题是，壳内、外带电体在相互影响下是否会达成另一种与此不同的平衡分布呢？唯一性定理告诉我们，这是不可能的。因为a、c两图中内部空间的边界条件相同（腔内表面电势为 $0$，内部带电体上总电量 $Q$ 给定），从而不管外部是否有带电体，内部的恒定分布是唯一的。这便是壳对内部的静电屏蔽效应。同理，因c、b两图中外部空间的边界条件相同（壳外表面电势为 $0$，外部带电体上总电量 $Q_1, Q_2, \ldots$ 给定，无穷远电势为 $0$），从而不管内部是否有带电体，外部的恒定分布是唯一的。这便是壳对外部的静电屏蔽效应。
\end{so}


\section{唯一性定理——从麦克斯韦方程出发}

上一节我们用物理思路给出了唯一性定理的证明，依靠的是对电场线和等势面的直观分析。现在，我们换一个角度——直接从麦克斯韦方程出发，用更数学化的语言再证明一遍。

\subsection[从麦克斯韦方程到泊松方程]{\itr{Maxwell's equations}{麦克斯韦方程}到\itr{Poisson equation}{泊松方程}}

在静电场中，麦克斯韦方程简化为：
\[\nabla \cdot \vec{E} = \frac{\rho}{\varepsilon_0}, \qquad \nabla \times \vec{E} = 0.\]
无旋性保证了我们可以引入电势 $U$，使得 $\vec{E} = -\nabla U$。代入散度方程，得到泊松方程：
\[\nabla^2 U = -\frac{\rho}{\varepsilon_0}.\]
在无电荷区域，退化为\itr{Laplace equation}{拉普拉斯方程} $\nabla^2 U = 0$。

对于存在导体的边值问题，边界条件通常有两种提法：
\begin{Itemize}
    \item \itr{Dirichlet}{狄利克雷}边界条件：给定边界上的电势值 $U|_S = U_0$；
    \item \itr{Neumann}{诺伊曼}边界条件：给定边界上的电势法向导数 $\left.\frac{\partial U}{\partial n}\right|_S$（等价于给定导体表面的电荷密度，因为 $\sigma_e = -\varepsilon_0 \frac{\partial U}{\partial n}$\mgnote{见\refleaftext{law8.3}。}）。
\end{Itemize}
唯一性定理断言：Dirichlet 边值问题的解是唯一的，Neumann 边值问题的解在相差一个常数意义下是唯一的（即电场唯一）。

\subsection[格林第一恒等式]{\itr{Green's first identity}{格林第一恒等式}}

证明的关键工具是格林第一恒等式。对于任意两个在体积 $V$ 及其边界 $S$ 上具有连续二阶导数的标量函数 $\phi$ 和 $\psi$，有：
\[\int_V \left( \phi \nabla^2 \psi + \nabla \phi \cdot \nabla \psi \right) \dif V = \oint_S \phi \frac{\partial \psi}{\partial n} \dif S.\]
读者暂时不需要纠结这个恒等式的推导——只需要理解为它把对空间的复杂体积分转化为了对边界的面积分。这非常关键，因为边值问题的核心矛盾恰恰在此：未知的是空间中的电势，已知的恰好是边界上的条件！这个等式可以由高斯散度定理直接推出：
\[\oint_S \phi \frac{\partial \psi}{\partial n} \dif S = \oint_S \phi \nabla \psi \cdot \dif \vec{S} = \int_V \nabla \cdot (\phi \nabla \psi) \dif V = \int_V \left( \nabla \phi \cdot \nabla \psi + \phi \nabla^2 \psi \right) \dif V.\]

\subsection{Dirichlet 边值问题的唯一性}

设在区域 $V$ 内有两个电势分布 $U_1$ 和 $U_2$，它们满足相同的泊松方程和相同的 Dirichlet 边界条件：
\[\nabla^2 U_1 = \nabla^2 U_2 = -\frac{\rho}{\varepsilon_0} \quad (\text{在 } V \text{ 内}), \qquad U_1|_S = U_2|_S = U_0 \quad (\text{在边界 } S \text{ 上}).\]
令 $W = U_1 - U_2$，则 $W$ 满足：
\[\nabla^2 W = 0 \quad (\text{在 } V \text{ 内}), \qquad W|_S = 0 \quad (\text{在边界 } S \text{ 上}).\]
在格林第一恒等式中取 $\phi = \psi = W$，得到：
\[\int_V \left( W \nabla^2 W + \nabla W \cdot \nabla W \right) \dif V = \oint_S W \frac{\partial W}{\partial n} \dif S.\]
由于 $\nabla^2 W = 0$ 且 $W|_S = 0$，上式简化为：
\[\int_V |\nabla W|^2 \dif V = 0.\]
被积函数 $|\nabla W|^2 \ge 0$，故在整个 $V$ 内必须有 $\nabla W = 0$，即 $W$ 为常数。又因为 $W$ 在边界 $S$ 上为 $0$，连续性要求 $W \equiv 0$ 处处成立。

从能量的视角看，这个结论格外自然：$\int_V |\nabla W|^2 \dif V$ 对应的正是静电场的能量。差解 $W$ 在边界上全为零——就像一个被"关掉了所有电源"的系统——又没有能量，那它里面自然不可能有任何电场波动，只能是处处为常数。对于 Dirichlet 问题，边界上 $W=0$，故 $W \equiv 0$，唯一性得证。对于 Neumann 问题，$W$ 为任意常数，电场 $\vec{E} = -\nabla W$ 仍然唯一。因此
\[U_1 = U_2, \quad \vec{E}_1 = \vec{E}_2.\]
Dirichlet 问题的解是\textbf{唯一}的。

\subsection{Neumann 边值问题的唯一性}

Neumann 问题给定的是边界上的法向导数。设 $U_1$ 和 $U_2$ 满足相同的泊松方程和相同的 Neumann 边界条件：
\[\nabla^2 U_1 = \nabla^2 U_2 = -\frac{\rho}{\varepsilon_0}, \qquad \left.\frac{\partial U_1}{\partial n}\right|_S = \left.\frac{\partial U_2}{\partial n}\right|_S.\]
同样令 $W = U_1 - U_2$，则：
\[\nabla^2 W = 0, \qquad \left.\frac{\partial W}{\partial n}\right|_S = 0.\]
再次使用格林第一恒等式（$\phi = \psi = W$）：
\[\int_V |\nabla W|^2 \dif V = \oint_S W \frac{\partial W}{\partial n} \dif S = 0.\]
同样得到 $\nabla W = 0$ 处处成立，即 $W$ 为常数。这意味着
\[U_1 = U_2 + \text{const}, \quad \vec{E}_1 = \vec{E}_2.\]
两个电势分布可以相差一个任意常数，但电场分布是\textbf{唯一}的。这个常数的不确定性正是电势零点选取自由度的体现——只要有一个导体的电势被确定，整个空间的电势分布就唯一了。


\section{电像法}

现在，我们终于可以回答第一节留下的那个问题了：\textbf{凭什么可以用像电荷替换导体？}

以第一节中最简单的例子——点电荷 $+q$ 放在无限大接地导体板前——来审视。我们"猜"出了一个解：在导体板后方对称位置放置一个 $-q$，然后假装导体板不存在。这个猜出来的电场分布，在导体板所在平面上恰好满足 $U=0$（因为 $+q$ 和 $-q$ 到平面上任意点的距离相等，电势相互抵消），在无穷远处也满足 $U \to 0$，在点电荷所在位置满足正确的奇性。所有的边界条件都与原问题一致。

而唯一性定理告诉我们：满足同一组边界条件的解只有一个。既然我们的"像电荷解"满足了全部边界条件，那它就是\textbf{唯一正确的解}。不需要再为导体表面那复杂的感应电荷分布而烦恼了。

这就是电像法的核心逻辑：
\begin{Itemize}
    \item \textbf{猜}：用假想的像电荷替换导体，猜测一个满足边界条件的简单解；
    \item \textbf{验}：验证该解在感兴趣的区域满足全部边界条件；
    \item \textbf{定}：由唯一性定理，该解就是问题的唯一正确答案。
\end{Itemize}


\section{接地导体平板前的点电荷}

有了唯一性定理的保障和电像法的核心思路，我们现在可以系统地处理几个经典的例子。首先是最简单的——接地导体平板。

回顾第一节中图4的 $U=0$ 等势面。设导体平板与点电荷相距为 $a$，点电荷带电量为 $q$，导体平板接地。由于 $q$ 的存在，接地导体平板上将产生感应电荷分布。我们在与导体平板表面对称的位置放置像电荷 $-q$，则由 $q$ 和像电荷 $-q$ 即可算出导体平板右侧空间的场强分布——原本复杂的感应电荷，被一个简单的点电荷所替代。

\begin{singlefigure}[接地导体平板前的点电荷]{electroimage_5.png}[0.55]
    以接地平板为 $z=0$ 界面建立坐标系，实电荷 $+q$ 位于轴上距离 $a$ 处，像电荷 $-q$ 位于对称位置 $-a$ 处。空间中一点 $P$ 坐标为 $(z,r)$。
\end{singlefigure}

由于系统的轴对称性，过对称轴取一平面建立 $Ozr$ 坐标，如图5所示。图中 $P$ 点坐标为 $(z, r)$，$P$ 点处的电场强度为
\[
    \vec{E}_P = -\frac{1}{4\pi\varepsilon_0} \frac{q}{r_2^2} \cdot \frac{\vec{r}_2}{r_2} + \frac{1}{4\pi\varepsilon_0} \frac{q}{r_1^2} \cdot \frac{\vec{r}_1}{r_1} = \frac{q}{4\pi\varepsilon_0} \left( \frac{\vec{r}_1}{r_1^3} - \frac{\vec{r}_2}{r_2^3} \right),
\]
其中
\[
    r_1^2 = r^2 + (z-a)^2, \quad r_2^2 = r^2 + (z+a)^2,
\]
\[
    \vec{r}_1 = (z-a)\hat{z} + r\hat{r}, \quad \vec{r}_2 = (z+a)\hat{z} + r\hat{r}.
\]

利用导体表面附近法向电场与电荷面密度的关系 $E = \sigma_e/\varepsilon_0$\mgnote{见\refleaftext{law8.3}。}，在导体表面 $z=0$ 处，有
\[
    \vec{E}(r) = -\frac{1}{4\pi\varepsilon_0} \frac{2aq}{(r^2 + a^2)^{3/2}} \hat{z},
\]
从而得到导体平板上的感应电荷面密度
\[
    \sigma_e(r) = -\frac{aq}{2\pi (r^2 + a^2)^{3/2}}.
\]

接地金属导体平板上感应电荷的总量为
\[
    q' = \int_0^{\infty} \sigma_e(r) \, 2\pi r \dif r = -\frac{aq}{2} \int_0^{\infty} \frac{\dif(r^2)}{(r^2 + a^2)^{3/2}} = -q.
\]
全部感应电荷恰好等于像电荷的电量——这绝非巧合，正是像电荷法自洽性的体现。

接地导体平板对点电荷 $q$ 的吸引力（即感应电荷对 $q$ 的作用力），可以通过像电荷直接算出：
\[
    F = \frac{1}{4\pi\varepsilon_0} \frac{q^2}{4a^2},
\]
方向指向平板。这个结果用直接计算感应电荷分布产生的电场来求，积分将十分繁琐——电像法的威力，由此可见一斑。


\section{接地导体球与等势面}

接下来讨论更复杂的几何形状——导体球。先从一个预备观察开始。

\subsection{不等量异号电荷的零电势面}

设有一对电荷，带电量分别为 $-q_1$ 和 $+q_2$（$q_1, q_2 > 0$，$q_1 \neq q_2$），相隔一段距离。考察电势为零的等势面，如图6所示。在此等势面上任一点 $P$，离两点电荷距离分别为 $r_1$ 和 $r_2$，则
\[
    -\frac{1}{4\pi\varepsilon_0} \cdot \frac{q_1}{r_1} + \frac{1}{4\pi\varepsilon_0} \cdot \frac{q_2}{r_2} = 0,
\]
化简得
\[
    \frac{r_1}{r_2} = \frac{q_1}{q_2} = \text{定值}.
\]

即零电势等势面上任一点到两定点距离之比为一常量。根据几何学结论，动点 $P$ 的轨迹在三维空间中是一个\textbf{球面}。这个结论是后续讨论的基础。

\begin{singlefigure}[不等量异号电荷的零电势面]{electroimage_6.png}[0.55]
    两个不等量异号电荷 $-q_1$ 和 $+q_2$ 的系统。电势为零的等势面为闭合曲面（球面），包围了电量较小的一侧。
\end{singlefigure}

如果将此等势面用接地薄金属片代替，根据静电屏蔽原理（唯一性定理），整个空间分成内外两部分后，场强分布不会发生变化。即使取走内部的电荷，也不会改变外部空间的电场分布。

\subsection{接地导体球前的点电荷}

有了上述准备，便可以解决接地导体球与球外点电荷的问题了。设接地导体球半径为 $R$，球外点电荷 $q$ 离球心距离为 $d$，如图7所示。

\begin{singlefigure}[接地导体球前的点电荷]{electroimage_7.png}[0.55]
    导体球半径为 $R$，接地。实电荷 $+q$ 位于轴线上距球心 $d$ 处。像电荷 $q'$ 位于球心到 $q$ 的线段上，距球心 $x = R^2/d$。
\end{singlefigure}

根据系统的轴对称性，像电荷必在对称轴上。设像电荷带电量为 $q'$，离球心距离为 $x$。考察对称轴线与球面相交的两点 $A$（近点电荷侧）和 $B$（远点电荷侧），写出电势为零的方程：
\[
    \frac{q'}{R-x} + \frac{q}{d-R} = 0, \qquad \frac{q'}{R+x} + \frac{q}{d+R} = 0.
\]

联立解得
\[
    q' = -\frac{R}{d}q, \qquad x = \frac{R^2}{d}.
\]

这便是接地导体球问题的像电荷。注意两个有趣的特点：\textbf{像电荷的绝对值小于原电荷}（因为 $R/d < 1$），且 \textbf{$x < R$}（像电荷位于球内）。当 $d \gg R$（点电荷很远）时，$q' \to 0$，$x \to 0$——导体球的感应越来越弱，像电荷缩向球心，物理上完全合理。

有了像电荷，就可以求出球外空间的场强分布、球外点电荷受到的作用力、以及球面上任意点的感应电荷面密度。此外，导体球表面的感应电荷总量也可以通过球心处的电势叠加求得：
\[
    \frac{1}{4\pi\varepsilon_0} \frac{q}{d} + \frac{1}{4\pi\varepsilon_0} \frac{Q}{R} = 0,
\]
解得 $Q = -\frac{R}{d}q$，恰好等于像电荷电量 $q'$。


掌握了接地导体球的基本解法后，可以通过叠加原理处理更一般的情形。核心思路是：将复杂问题分解为"接地导体球 + 点电荷"（已解决）与"均匀带电导体球"（已知解）的叠加。

\subsection{不带电的孤立导体球}

若一个半径为 $R$ 的导体球不带电也不接地，球外距球心 $d$ 处有点电荷 $q$。求解球外空间的场强分布以及球外点电荷受到的作用力。

我们分两步构造解。\textbf{第一步}，假设导体球接地。由前面的结论，球面上感应电荷总量为 $q' = -\frac{R}{d}q$，这些感应电荷对球外电场的作用可用位于 $x = R^2/d$ 处的像电荷 $q'$ 等效。此时将接地导线切断，导体球带电量即为 $q'$，但仍保持等势。

\textbf{第二步}，现在导体球带电 $q'$，而问题要求球不带电（总电量为零）。怎么办呢？我们需要在球体上再加入 $-q'$ 的电量，同时必须保证球体仍为等势体——否则唯一性定理的前提就破坏了。唯一可行的方法是：\textbf{将 $-q'$ 均匀分布于球面}。根据导体静电平衡的性质，均匀带电球面在球外产生的电场，等价于将全部电荷集中于球心的点电荷。

因此，球外空间的场强由三个点电荷共同激发：
\begin{Itemize}
    \item 球外点电荷 $q$（位置：距球心 $d$）；
    \item 像电荷 $q' = -\frac{R}{d}q$（位置：距球心 $x = R^2/d$）；
    \item 球心处补偿电荷 $-q' = \frac{R}{d}q$。
\end{Itemize}

由唯一性定理，这个叠加解是问题所要求的唯一解。球外点电荷 $q$ 受到导体球的作用力，就是像电荷 $q'$ 和球心处 $-q'$ 对 $q$ 的库仑力之和。

\subsection{带总电量 $Q$ 的导体球}

若导体球带总电量 $Q$ 且不接地，球外距球心 $d$ 处有点电荷 $q$。这个情形与上面"不带电球"的处理几乎一样——唯一的区别在于球体的总电量不同。

沿用上述分步构造法：接地导体球上的感应电荷 $q' = -\frac{R}{d}q$ 仍由像电荷 $q'$ 等效；但此时球体需要的总电量为 $Q$（而非零），所以球心处应放置的补偿电荷为 $Q - q'$（而非 $-q'$）。因此：
\begin{Itemize}
    \item 像电荷 $q' = -\frac{R}{d}q$，位于距球心 $x = R^2/d$ 处；
    \item 球心处放置补偿电荷 $Q - q' = Q + \frac{R}{d}q$。
\end{Itemize}
两个像电荷共同决定了球外电场和球外点电荷的受力。

\subsection{电势固定为 $U$ 的导体球}

若导体球维持在恒定电势 $U$（例如通过电池与地相连），球外距球心 $d$ 处有点电荷 $q$。同样用叠加法处理。

接地导体球的像电荷仍为 $q' = -\frac{R}{d}q$（位置 $x = R^2/d$）。此外，为维持球体电势为 $U$，球心处需放置一个合适的电荷 $Q_U$。这个 $Q_U$ 可以利用孤立导体球的电容公式求得：半径为 $R$ 的孤立导体球，电容为
\[
    C = 4\pi\varepsilon_0 R,
\]
带上电量 $Q_U$ 后，其电势为
\[
    U = \frac{1}{4\pi\varepsilon_0} \frac{Q_U}{R},
\]
由此得 $Q_U = 4\pi\varepsilon_0 R U$。因此，球外电场由像电荷 $q'$ 和球心处电荷 $Q_U = 4\pi\varepsilon_0 R U$ 共同决定。

可以看到，这三种情形的处理遵循完全相同的逻辑：\textbf{接地球的像电荷 $q'$ 负责满足"球面为零等势面"，球心处的点电荷负责满足"总电量"或"电势"的边界条件}。有了唯一性定理做担保，叠加法给出的就是唯一正确答案。

然而，以上三种情形的叠加法之所以简便，是因为它们共享一个前提：每个被叠加的"子问题"（接地导体球 $+$ 点电荷、孤立均匀带电球）本身已经有现成的解。当导体的几何配置变得更复杂——比如下面这个例子——没有现成的子问题可以叠加，此时就需要用到电像法更强大的一个层面：\textbf{迭代}。

\begin{ex}[两个接触导体球的电容]
    两个半径均为 $R$ 的导体球相互接触，形成一孤立导体。试利用电像法求此孤立导体的电容。
\end{ex}

\textbf{分析}：孤立导体的电容定义为 $C = Q/U_0$，其中 $U_0$ 是电势，$Q$ 是达到此电势所需的总电量。若孤立导体外形很对称（例如单个球体），充电后表面电荷分布均匀\mgnote{均匀分布时，球外电场等价于全部电荷集中于球心的点电荷。}，则电场强度分布易求，$U_0$ 与 $Q$ 的关系一举可知，问题得解。

然而本题的孤立导体是两个相互接触的等径球体，不具有球对称性——充电后表面电荷分布\textbf{不均匀且未知}。这就卡住了：不知道电荷怎么分布，就没法求外电场；不知道外电场，就没法建立 $Q$ 与 $U_0$ 的关系。常规的叠加法在这里无计可施。

电像法的核心洞察是：\textbf{将两球面上复杂的不均匀电荷分布，用一系列位于球内的虚设点电荷来等效代替}。这些点电荷在球外激发的电场，与真实的面电荷分布在球外激发的电场完全一致。一旦这些虚设电荷在两球面上产生的电势都等于 $U_0$，那么这些虚设电荷的\textbf{总电量}就是使孤立导体达到电势 $U_0$ 所需的实际电量 $Q$。

\begin{so}[两个接触导体球的电容]
    \begin{singlefigure}[两个接触导体球的镜像法迭代]{electroimage_10.png}[0.7]
        两球心处虚设第一对点电荷 $q_1$，随后在球内逐次引入 $q_2, q_3, \ldots$ 及其对应距离参数 $b_2, b_3, \ldots$
    \end{singlefigure}

    设孤立导体的电势为 $U_0$。\textbf{第一步}，在两个球心处各虚设一个点电荷 $q_1 = 4\pi\varepsilon_0 R U_0$。这样，每个 $q_1$ \textbf{在自己的球面上}产生的电势恰好为 $U_0$——乍看之下，目标已经达成。

    但仔细一想便发现漏洞：$q_1$ 不只是在自己的球面上产生电势，它同时也在对方的球面上产生附加电势。以左球心的 $q_1$ 为例，它在右球面上各点的电势\textbf{并不是常数}（因为距离不同），而且整体抬升了右球面的电势——结果就是两个球面的实际电势都\textbf{大于} $U_0$，且不是等势面。第一步的猜测只是一个粗糙的初始近似。

    \textbf{第二步}：为消除 $q_1$ 在对方球面上产生的附加电势，需在对方球内适当位置引入第二对虚设点电荷 $q_2$。此处直接引用接地导体球的像电荷公式（§6.2）：对面球心处的 $q_1$（视为外部点电荷，距本地球心 $2R$）在本地球内的像电荷为
    \[
        q_2 = -\frac{R}{2R}\, q_1 = -\frac{q_1}{2}, \qquad
        b_2 = \frac{R^2}{2R} = \frac{R}{2},
    \]
    其中 $b_2$ 是 $q_2$ 与所在球心之间的距离。这样，左球心处的 $q_1$ 与右球内 $b_2$ 处的 $q_2$，联合作用使右球面成为等势面（电势为零——但球面不是零电势，后续迭代会校正总电势）。

    然而，新引入的 $q_2$ 又会在\textbf{原}球面上产生附加电势！为了消除这一轮附加电势，需再次利用像电荷公式：

    \textbf{第三步}：$q_2$（位于对面球内，距本地球心 $2R - b_2$）在本地球内的像电荷为
    \[
        q_3 = -\frac{R}{2R - b_2}\, q_2 = \frac{q_1}{3}, \qquad
        b_3 = \frac{R^2}{2R - b_2} = \frac{2}{3}R.
    \]

    如此继续——每引入一对新的虚设电荷消除上一轮的附加电势，自身又产生新的附加电势——永无止境。这就是电像法的\textbf{迭代}机制。

    \textbf{第 $n$ 对虚设电荷的递推}：由上述逻辑，第 $n$ 对虚设点电荷的电量 $q_n$ 及其与所在球心的距离 $b_n$ 满足
    \begin{align}
        q_n & = -\frac{R}{2R - b_{n-1}}\, q_{n-1}, \\
        b_n & = \frac{R^2}{2R - b_{n-1}}.
    \end{align}

    两式相除，得简洁的关系 $b_n = -R q_n / q_{n-1}$。由此消去 $b_n$、$b_{n-1}$，递推式化简为
    \[
        \frac{1}{q_n} + \frac{2}{q_{n-1}} + \frac{2}{q_{n-2}} = 0.
    \]

    代入初值 $q_1 = 4\pi\varepsilon_0 R U_0$，$q_2 = -q_1/2$，递推得
    \[
        q_3 = \frac{q_1}{3},\quad q_4 = -\frac{q_1}{4},\quad \ldots,\quad
        q_n = (-1)^{n+1}\,\frac{q_1}{n}.
    \]

    验证可知此通解满足递推关系。注意 $q_n$ 的符号交替变化，且 $|q_n| \propto 1/n$ 缓慢衰减——这意味着收敛需要无穷多项，但级数确实收敛（调和级数交替项）。

    \textbf{从迭代到总电容}：当 $n \to \infty$ 时，两球面的电势精确达到 $U_0$，所有虚设电荷的总电量即为孤立导体的实际带电量：
    \[
        Q = \sum_{n=1}^{\infty} 2q_n
        = 2q_1 \left( 1 - \frac{1}{2} + \frac{1}{3} - \frac{1}{4} + \cdots \right).
    \]

    括号中的交错调和级数正是 $\ln 2$ 的级数展开：
    \[
        \ln 2 = 1 - \frac{1}{2} + \frac{1}{3} - \frac{1}{4} + \cdots,
    \]

    因此
    \[
        Q = 2q_1 \ln 2 = 8\pi\varepsilon_0 R U_0 \ln 2,
        \qquad
        \boxed{C = \frac{Q}{U_0} = 8\pi\varepsilon_0 R \ln 2}.
    \]

    \textbf{回顾}：这个例子完美展示了电像法并非要求一步到位，而可以一次次运用电像法进行对于实际情况的迭代逼近。每次迭代都消除了上一轮的附加电势，最终在无穷多轮迭代后，两个球面上的电势精确达到 $U_0$，从而求得了孤立导体的总电量和电容。
\end{so}


\section{循迹之边界——电像法能走多远？}

至此，我们用电像法处理了平板、球等几何形状。读者自然会问：电像法到底能解多少种形状？有没有一个"可解集"？

\subsection{经典可解形状}

历史上，人们陆续发现以下几类导体边界存在简单的点像电荷解：

\begin{Itemize}
    \item \textbf{无限大平面}：点电荷的像是对称位置的等量异号电荷（§5）；
    \item \textbf{导体球}：点电荷的像位于球内，$q' = -\frac{R}{d}q$，$x = \frac{R^2}{d}$（§6）；
    \item \textbf{无限长导体圆柱}：线电荷的像满足 $\lambda' = -\lambda$，$x = \frac{R^2}{a}$（习题2）；
    \item \textbf{导体劈}：夹角为 $\pi/n$（$n$ 为整数）的两个相交接地导体平面之间的区域。例如直角劈（$90^\circ$）需要 3 个像电荷，$60^\circ$ 劈需要 5 个——本质上是对平面镜像的多次反射。
\end{Itemize}

这些形状的共同特征是：它们的等势面恰好与一组点电荷（或线电荷）系统的某个等势面重合。换句话说，电像法成功的条件是——\textbf{导体的边界必须是某个已知电荷系统的等势面}。

那么，超出这个"经典可解集"呢？对于椭球面、双曲面、抛物面这些更一般的二次曲面，电像法是否仍然有效？

\subsection{二次曲面的统一格林函数解}

从数学上看，旋转二次曲面可以统一用极坐标方程描述：
\[
    \rho = \frac{p'}{1 + e\cos\theta},
\]
其中 $e$ 为离心率，$p'$ 为焦点参数。四种二次曲线的区别仅在于 $e$ 的取值：
\[
    \begin{aligned}
        \text{椭圆 } (0 \le e < 1), \quad
        \text{抛物线 } (e = 1), \quad
        \text{双曲线 } (e > 1), \quad
        \text{圆 } (e = 0).
    \end{aligned}
\]

在\textbf{近轴条件}（点电荷位于对称轴上，且考察区域限于轴附近）下，像电荷的位置和电量可以写成\textbf{统一形式}：
\[
    \boxed{ x = \frac{p'^{\,2}}{d\,(1+e)^2}, \qquad
        q = \frac{p'}{d\,(1+e)}\, q_{\text{实}} }.
\]

这一公式同时涵盖了四种经典二次曲面！代入各自参数后，得到：

\begin{center}
    \begin{tabular}{c|c|c}
        \toprule
        \textbf{曲面类型}        & \textbf{像电荷位置 $x$}       & \textbf{像电荷电量 $q$}                      \\
        \midrule
        球面（半径 $R$）           & $x = \dfrac{R^2}{d}$     & $q = \dfrac{R}{d}\,q_{\text{实}}$        \\[6pt]
        椭球面（长半轴 $a$，短半轴 $b$） & $x = \dfrac{(a-c)^2}{d}$ & $q = \dfrac{b^2}{d(a+c)}\,q_{\text{实}}$ \\[6pt]
        双曲面（实半轴 $a$）         & $x = \dfrac{(c-a)^2}{d}$ & $q = \dfrac{c-a}{d}\,q_{\text{实}}$      \\[6pt]
        抛物面（焦点参数 $p$）        & $x = \dfrac{p^2}{4d}$    & $q = \dfrac{p}{2d}\,q_{\text{实}}$       \\
        \bottomrule
    \end{tabular}
\end{center}
其中 $c = \sqrt{a^2 - b^2}$（椭球）或 $c = \sqrt{a^2 + b^2}$（双曲面），$d$ 为实电荷到坐标原点的距离。

球面的结果（第一行）正是我们熟悉的 $x = R^2/d$ 和 $q' = -\frac{R}{d}q$（符号差异来自 $q$ 的定义方向）。这验证了统一公式的自洽性。

\subsection{近轴之外}

近轴近似固然简洁，但一旦离开对称轴，单个点像电荷不再能精确满足边界条件。远轴情形下的严格求解需要借助\itr{separation of variables}{分离变量法}（在共焦椭球坐标或双曲坐标下展开）或\itr{conformal mapping}{共形映射}等更系统的数学工具。此时"像"可能不再是离散的点电荷，而是分布于焦线上的连续电荷分布。其中所使用的数学物理方法非常复杂，已经远远超出了笔者的能力范围，并且脱离了电像法用简单直观的构造解决复杂问题的本质，因此这里仅可为读者简述。


\section{总结}

走到本章的终点，我们不妨回过头来，重新审视那个最初的问题：\textbf{电像法到底是什么？}

电像法诞生于 1847 年，由开尔文勋爵（William Thomson）在求解导体球前点电荷问题时提出，他将其命名为\itr{method of images}{电像法}。这个名字中的"像"（image）携带着双重的含义：\textbf{几何的镜像}——像电荷如同实电荷在导体表面反射出的虚影；\textbf{物理的镜像}——导体表面的感应电荷分布本身就是对外部电荷的"响应"，导体"看见"了外部电荷，调整自己表面的电荷分布以维持等势，而像电荷不过是把这种复杂的分布式响应浓缩成了几个简单的点。

这，便是电像法的第一个面向：\textbf{用简单替代复杂}。

但替代不能是空中楼阁。唯一性定理为这个替代提供了严格的保障——边界条件一旦给定，解便是唯一的。点电荷与像电荷的简单系统，和导体与点电荷的复杂系统之间，由此建立了一一对应的关系，完成了严谨的逻辑闭环。这，是第二个面向：\textbf{用严谨为直觉护航}。

而最有意思的是第三个面向：像电荷的位置和电量应该如何被构造出来？答案是一个字——"猜"。但这不是盲目的试探：平板前的 $-$q 来自反射对称性，球内的 $q'$ 来自反演对称性，接触球间的无穷级数来自迭代自相似性。每一个成功的猜测背后，本质都是对几何结构深层对称性的洞察。这便是电像法的第三个面向——\textbf{用对称性指引直觉}。

如此，便是为什么笔者将电像法独立出来作为一个循迹之外章节的原因：它不仅仅是一个解题技巧，更是一种思维方式。希望读者在合上这份讲义之后，记住的不是那几个公式，而是这样一种习惯——面对一个看似无路可走的物理问题时，先环顾四周，寻找属于它的那面"镜子"。

\let\thesection\oldthesection
\let\thesubsection\oldthesubsection
